% =====================================================================
%  A COUNTING BOUND ON TRANSVERSE RELAXATION
%  Schmidt-Rank Growth, the Margolus--Levitin Ceiling, and a
%  Perturbation-Independent Loschmidt-Echo Floor
%
%  Self-contained. No external framework assumed.
% =====================================================================
\documentclass[twocolumn,10pt,a4paper]{article}

\usepackage[utf8]{inputenc}
\usepackage[T1]{fontenc}
\usepackage{lmodern}
\usepackage[a4paper,top=2cm,bottom=2.2cm,left=1.8cm,right=1.8cm,columnsep=0.6cm]{geometry}
\usepackage{amsmath,amssymb,amsthm,mathtools}
\usepackage{bm}
\usepackage{booktabs}
\usepackage{array}
\usepackage{enumitem}
\usepackage{microtype}
\usepackage{xcolor}
\usepackage[round,sort&compress]{natbib}
\usepackage[colorlinks=true,linkcolor=blue!55!black,
            citecolor=blue!55!black,urlcolor=blue!55!black]{hyperref}
\usepackage{fancyhdr}

\pagestyle{fancy}
\fancyhf{}
\fancyhead[LE,RO]{\thepage}
\fancyhead[RE]{\small A Counting Bound on Transverse Relaxation}
\fancyhead[LO]{\small}
\renewcommand{\headrulewidth}{0.4pt}

\newtheoremstyle{thmstyle}{\topsep}{\topsep}{\itshape}{}{\bfseries}{.}{.5em}{}
\theoremstyle{thmstyle}
\newtheorem{theorem}{Theorem}[section]
\newtheorem{lemma}[theorem]{Lemma}
\newtheorem{proposition}[theorem]{Proposition}
\newtheorem{corollary}[theorem]{Corollary}
\theoremstyle{definition}
\newtheorem{definition}[theorem]{Definition}
\newtheorem{assumption}{Assumption}
\newtheorem{example}[theorem]{Example}
\theoremstyle{remark}
\newtheorem{remark}[theorem]{Remark}

\newcommand{\R}{\mathbb{R}}
\newcommand{\N}{\mathbb{N}}
\newcommand{\C}{\mathbb{C}}
\newcommand{\kB}{k_{\mathrm{B}}}
\newcommand{\Hilb}{\mathcal{H}}
\newcommand{\Tr}{\operatorname{Tr}}
\newcommand{\rank}{\operatorname{rank}}
\newcommand{\sysrho}{\rho_{\mathrm{S}}}
\newcommand{\Svn}{S_{\mathrm{vN}}}
\newcommand{\Meff}{M_{\mathrm{eff}}}
\newcommand{\Tone}{T_{1}}
\newcommand{\Ttwo}{T_{2}}
\newcommand{\Tstar}{T_{2}^{\ast}}
\newcommand{\dd}{\mathrm{d}}
\newcommand{\Ham}{\mathcal{H}}
\DeclareMathOperator{\Var}{Var}

% =====================================================================
\title{\textbf{A Counting Bound on Transverse Relaxation:\\[2pt]
Schmidt-Rank Growth, the Margolus--Levitin Ceiling,\\
and a Perturbation-Independent Loschmidt-Echo Floor}}

\author{}
\date{\today}
% =====================================================================

\begin{document}
\maketitle

% =====================================================================
\begin{abstract}
\noindent
The transverse relaxation time $\Ttwo$ of a spin ensemble, and the
irreducible attenuation of its Hahn and Loschmidt echoes, are normally
computed from a microscopic relaxation theory (Bloch--Redfield, or its
strong-coupling Anderson--Kubo limit) in which the relaxation rate is a
second-moment integral of a bath correlation function. We present an
independent, information-theoretic derivation that does not begin from a
correlation function but from a single counting principle: the number of
mutually distinguishable states a system can resolve under a Hamiltonian
of energy spread $\Delta E$ grows no faster than the Margolus--Levitin
ceiling $2\Delta E/\pi\hbar$. We make this concrete by identifying the
``count'' with the effective Schmidt rank $\Meff(t)=\exp \Svn(\sysrho(t))$
of the reduced spin state, where $\Svn$ is the von Neumann entropy of the
spin density operator after the bath is traced out. Three results follow.
\textbf{(i)} A \emph{counting bound on relaxation}: the transverse
relaxation rate is bounded by the rate of Schmidt-rank growth, which is
in turn bounded by the Margolus--Levitin ceiling evaluated at the
spin--bath coupling energy; saturation of the bound reproduces the
Anderson--Kubo rigid-lattice rate $1/\Ttwo\simeq\langle\delta\omega^{2}
\rangle^{1/2}$ with no fitted constant. \textbf{(ii)} A \emph{derivation,
not a postulate, of monotonicity}: starting from a strictly
time-reversal-invariant total Hamiltonian, the reduced spin entropy
$\Svn(\sysrho(t))$ is non-decreasing under coarse-grained evolution by
the data-processing inequality, so the count is monotone although the
microscopic dynamics is reversible---the arrow enters solely through the
partial trace, and the spin echo separates the reversible
(product-state) part of dephasing from the irreversible
(entanglement-committed) part exactly. \textbf{(iii)} A
\emph{perturbation-independent echo floor}: counting the distinguishable
many-body trajectories accessible after $n$ coherent steps over $d$
coupling channels gives $T(n,d)=d(d+1)^{n-1}$, whose logarithmic growth
rate $\ln(d{+}1)$ per step is independent of the coupling strength;
converted to a physical rate this predicts a Loschmidt-echo decay floor
$1/\tau_{\mathrm{floor}}=(\omega_{\mathrm{loc}}/2\pi)\ln(d{+}1)$ that does
not recede as the reversal engineering improves. The prediction
reproduces the observed environment-independence of polarization-echo
decay and, unlike the semiclassical Lyapunov account, predicts a nonzero
floor for non-chaotic and integrable many-body spin systems---a
falsifiable point of departure. We state the assumptions explicitly,
prove each result, and identify the experiments that would confirm or
refute the framework.
\end{abstract}

\noindent\textbf{Keywords:} transverse relaxation; spin echo; Loschmidt
echo; Schmidt rank; von Neumann entropy; Margolus--Levitin bound; quantum
speed limit; data-processing inequality; environment-independent
decoherence; irreversibility.

\tableofcontents

% =====================================================================
\section{Introduction}
\label{sec:intro}
% =====================================================================

\subsection{The question}

A precessing spin ensemble in a static magnetic field loses transverse
coherence on two timescales \citep{bloch1946,abragam1961,slichter1990}.
The faster, $\Tstar$, reflects static spread in precession frequencies
across the sample; it is reversible, and the Hahn echo
\citep{hahn1950} undoes it exactly with a single refocusing pulse. The
slower, $\Ttwo$, reflects time-dependent fluctuating fields; it is not
undone by refocusing, and it sets the floor of the echo amplitude. The
standard account of $\Ttwo$ is Bloch--Redfield theory
\citep{redfield1957,abragam1961}: the relaxation rate is a spectral
integral of a bath autocorrelation function, reducing to
$1/\Ttwo=\tfrac12\gamma^{2}\langle\delta B^{2}\rangle\tau_{c}$ in the
motional-narrowing limit and to the rigid-lattice second moment
$1/\Ttwo\simeq\langle\delta\omega^{2}\rangle^{1/2}$ when the bath is slow
\citep{anderson1953,kubo1962}.

This theory is quantitatively successful and we do not dispute it. We
ask a different question. The relaxation rate, in the standard account,
is the output of a correlation function whose form is supplied by the
microscopic model. Is there a derivation of the \emph{same number} from a
principle that does not presuppose the correlation function---a principle
about how fast distinguishable information can be generated at all? And
does that principle say anything the correlation-function account does
not, in the regime where the two might diverge?

We answer both questions affirmatively. The organizing principle is a
\emph{quantum speed limit}: the Margolus--Levitin theorem
\citep{margolus1998,mandelstam1945,deffner2017}, which bounds the rate at
which a system can pass between distinguishable states by its energy
spread. We convert this into a bound on a \emph{counting} quantity---the
effective number of states the environment has resolved about the
spins---and show that the relaxation rate is the saturation of this
bound.

\subsection{The three layers of the argument}

The paper has a deliberate three-layer structure, each layer answering a
distinct objection one might raise against the previous one.

\paragraph{Foundation (Section~\ref{sec:ceiling}): a counting ceiling
fixes the rate constant.} We define a count---the effective Schmidt rank
of the reduced spin state---and prove that its rate of growth is bounded
by the Margolus--Levitin ceiling evaluated at the spin--bath interaction
energy. Saturating the ceiling reproduces the Anderson--Kubo rigid-lattice
relaxation rate with no adjustable proportionality constant. This answers
the objection ``the rate constant is fitted, not derived.''

\paragraph{Mechanism (Section~\ref{sec:mechanism}): monotonicity is
derived, not assumed.} We then show \emph{why} the count grows at all,
and only forward. Starting from a strictly time-reversal-invariant total
Hamiltonian for spins-plus-bath, the reduced spin entropy is
non-decreasing under the coarse-grained (bath-averaged) dynamics by the
data-processing inequality for completely positive trace-preserving maps.
The microscopic dynamics is reversible; the count is monotone; the only
step that breaks the symmetry is the partial trace. The Hahn echo is the
experimentum crucis: it separates the part of dephasing that leaves the
spin--bath state a product (reversible, refocusable) from the part that
entangles spins with the bath (irreversible, committed). This answers the
objection ``the monotonicity is true by definition, not earned from
reversible mechanics.''

\paragraph{Headline (Section~\ref{sec:floor}): a
perturbation-independent floor.} Finally we count the distinguishable
many-body trajectories a spin system can traverse in $n$ coherent steps
when it is coupled through $d$ channels. The count is $T(n,d)=d(d+1)^{n-1}$,
and its logarithmic growth rate per step, $\ln(d+1)$, depends only on the
channel multiplicity $d$ and \emph{not} on the coupling strength.
Converted to a physical rate through the local precession frequency, this
predicts an irreducible Loschmidt-echo decay rate
$1/\tau_{\mathrm{floor}}=(\omega_{\mathrm{loc}}/2\pi)\ln(d+1)$ that does
not improve as the experimental reversal is made more perfect. This
reproduces the central empirical finding of polarization-echo experiments
\citep{levstein1998,pastawski2000,jalabert2001}---decay independent of the
residual imperfection---and makes a prediction the semiclassical Lyapunov
account does not: a nonzero floor even for integrable, non-chaotic spin
systems.

\subsection{Relation to existing theory}

We are not proposing a replacement for relaxation theory; we are
proposing a bound that the standard theory must respect and that, at
saturation, the standard theory realizes. The relationship is the same as
that between the Carnot bound and a specific heat engine: the bound is
derived from a general principle (here, the quantum speed limit), and a
particular microscopic model either saturates it or falls below it. The
content of the paper is (a) that transverse relaxation \emph{saturates}
the counting ceiling in the rigid-lattice regime, (b) that the
monotonicity of the count is forced by open-system structure rather than
assumed, and (c) that one consequence---the echo floor---is both
quantitative and, in a specific regime, distinguishable from the
orthodox semiclassical prediction.

% =====================================================================
\section{Setup and assumptions}
\label{sec:setup}
% =====================================================================

We fix notation and state every assumption explicitly, so that each later
theorem can be traced to the hypotheses it requires.

\subsection{The spin--bath system}

\begin{definition}[Spin--bath decomposition]
\label{def:sysbath}
The total Hilbert space is $\Hilb=\Hilb_{\mathrm{S}}\otimes
\Hilb_{\mathrm{B}}$, where $\Hilb_{\mathrm{S}}$ carries the observed spin
ensemble ($N$ spins, $\dim\Hilb_{\mathrm{S}}=2^{N}$) and
$\Hilb_{\mathrm{B}}$ carries everything else with which the spins
interact (other spins, lattice phonons, molecular degrees of freedom). The
total Hamiltonian is
\begin{equation}
H = H_{\mathrm{S}}\otimes\mathbb{1}
  + \mathbb{1}\otimes H_{\mathrm{B}}
  + H_{\mathrm{SB}},
\label{eq:Htotal}
\end{equation}
with $H_{\mathrm{SB}}$ the spin--bath coupling. The reduced spin state is
$\sysrho(t)=\Tr_{\mathrm{B}}\,\rho(t)$, where $\rho(t)$ is the state of
the closed total system.
\end{definition}

\begin{assumption}[Reversible substrate]
\label{ass:reversible}
The total Hamiltonian \eqref{eq:Htotal} is time-independent and Hermitian,
and is time-reversal invariant: there is an antiunitary $\Theta$ with
$\Theta H\Theta^{-1}=H$. The total evolution $U(t)=e^{-iHt/\hbar}$ is
unitary, hence exactly reversible.
\end{assumption}

This is the crux of Loschmidt's objection \citep{loschmidt1876}: the
substrate is reversible, and any irreversibility we derive must therefore
not be smuggled in through a non-reversible postulate. Assumption~%
\ref{ass:reversible} holds it fixed.

\begin{assumption}[Initially uncorrelated preparation]
\label{ass:product}
At $t=0$ the spin ensemble is prepared in a definite state uncorrelated
with the bath: $\rho(0)=\sysrho(0)\otimes\rho_{\mathrm{B}}(0)$. In the
spin-echo context this is the post-$90^{\circ}$-pulse transverse state.
\end{assumption}

Assumption~\ref{ass:product} is the standard and physically warranted
initial condition for a relaxation experiment: the pulse prepares the
spins independently of the bath microstate. It is the only place an
``initial condition'' enters, and it is the weakest possible one (a
product state, not a low-entropy fine-tuning).

\begin{assumption}[Coupling energy scale]
\label{ass:coupling}
The spin--bath coupling has a finite energy spread in the relevant state,
\begin{equation}
\Delta E_{\mathrm{SB}}
  := \bigl(\langle H_{\mathrm{SB}}^{2}\rangle
          - \langle H_{\mathrm{SB}}\rangle^{2}\bigr)^{1/2}
  = \hbar\,\langle\delta\omega^{2}\rangle^{1/2},
\label{eq:dE}
\end{equation}
where $\langle\delta\omega^{2}\rangle$ is the second moment of the
fluctuating local precession frequency. We write
$b:=\langle\delta\omega^{2}\rangle^{1/2}$ for this coupling frequency; it
is the single material parameter of the theory and is, in any concrete
system, computable from the dipolar (or hyperfine, or exchange) couplings
by the standard van Vleck moment formulae \citep{abragam1961}.
\end{assumption}

\subsection{The count}

The central object is not the cycle count of a clock but a count of
\emph{distinguishable spin states the bath has come to encode}.

\begin{definition}[Effective Schmidt rank / partition count]
\label{def:count}
Let $\sysrho(t)$ have von Neumann entropy
$\Svn(\sysrho)=-\Tr\!\bigl(\sysrho\ln\sysrho\bigr)$. The
\emph{effective Schmidt rank} of the spin--bath state, equivalently the
\emph{partition count} of the spin ensemble, is
\begin{equation}
\Meff(t) := \exp\!\bigl[\Svn(\sysrho(t))\bigr].
\label{eq:Meff}
\end{equation}
$\Meff$ is the effective number of orthogonal bath states correlated with
distinct spin states (the participation number of the Schmidt
decomposition of the global pure state across the S$\mid$B cut)
\citep{nielsen2010}.
\end{definition}

Three remarks fix the interpretation. First, $\Meff(0)=1$ by
Assumption~\ref{ass:product} (a product state has Schmidt rank one).
Second, $\Meff$ is dimensionless and basis-independent. Third, $\Meff$ is
exactly the quantity that controls echo attenuation: an echo recovers
coherence only to the extent that the spin state can be disentangled from
the bath by an operation on the spins alone, and the obstruction to that
is precisely $\Svn(\sysrho)>0$. We make this precise in
Section~\ref{sec:mechanism}.

\begin{remark}[Why a count, and why this count]
\label{rem:why-count}
The word ``count'' is used advisedly. $\ln\Meff=\Svn(\sysrho)$ is the
number of bits the bath has acquired about the spins; $\Meff$ is the
number of distinguishable alternatives those bits resolve. Relaxation,
decoherence, and echo attenuation are all, in this reading, the growth of
$\Meff$. The remainder of the paper bounds the \emph{rate} of that
growth, proves it \emph{monotone}, and computes its \emph{floor}.
\end{remark}

% =====================================================================
\section{Foundation: the counting ceiling on the relaxation rate}
\label{sec:ceiling}
% =====================================================================

We first bound how fast $\Meff$ can grow, using only the energy spread of
the dynamics, and show that saturating the bound returns the known
relaxation rate.

\subsection{The Margolus--Levitin ceiling}

\begin{lemma}[Distinguishability rate ceiling]
\label{lem:ML}
Let a quantum system evolve under a time-independent Hamiltonian of energy
spread $\Delta E$ (with the ground-state energy set to zero, mean energy
$\langle E\rangle$). The minimum time to evolve into an orthogonal
(perfectly distinguishable) state satisfies the Mandelstam--Tamm and
Margolus--Levitin bounds
\begin{equation}
\tau_{\perp}\;\geq\;\max\!\left\{\frac{\pi\hbar}{2\,\Delta E},\;
\frac{\pi\hbar}{2\,\langle E\rangle}\right\}.
\label{eq:ML}
\end{equation}
Consequently the rate at which mutually orthogonal states can be
generated is bounded:
\begin{equation}
\frac{\dd N_{\perp}}{\dd t}\;\leq\;\frac{2\,\Delta E}{\pi\hbar}.
\label{eq:rate-ceiling}
\end{equation}
\end{lemma}

\begin{proof}
The Mandelstam--Tamm bound $\tau_{\perp}\geq\pi\hbar/2\Delta E$ follows
from the energy--time uncertainty relation applied to the survival
amplitude $\langle\psi(0)|\psi(t)\rangle$
\citep{mandelstam1945,deffner2017}; the Margolus--Levitin bound
$\tau_{\perp}\geq\pi\hbar/2\langle E\rangle$ is an independent inequality
from the spectral decomposition of the survival amplitude
\citep{margolus1998}. Both are tight. A sequence of $N_{\perp}$ mutually
orthogonal states must be separated by at least $\tau_{\perp}$ in time, so
in elapsed time $t$ at most $t/\tau_{\perp}$ such states are reachable,
giving \eqref{eq:rate-ceiling}.
\end{proof}

\subsection{From the ceiling to the relaxation rate}

We now apply Lemma~\ref{lem:ML} to the bath's acquisition of which-spin
information. The relevant ``system'' for the speed limit is the
\emph{conditional bath state}: when the spins are in configuration $s$,
the bath evolves under an effective Hamiltonian
$H_{\mathrm{B}}^{(s)}=H_{\mathrm{B}}+\langle s|H_{\mathrm{SB}}|s\rangle$.
Two distinct spin configurations $s\neq s'$ drive the bath along
trajectories whose conditional Hamiltonians differ by
$\Delta H_{\mathrm{B}}=\langle s|H_{\mathrm{SB}}|s\rangle
-\langle s'|H_{\mathrm{SB}}|s'\rangle$, of energy spread of order
$\Delta E_{\mathrm{SB}}$ (Assumption~\ref{ass:coupling}).

\begin{theorem}[Counting bound on transverse relaxation]
\label{thm:counting-bound}
Under Assumptions~\ref{ass:reversible}--\ref{ass:coupling}, the rate of
growth of the partition count is bounded by the Margolus--Levitin ceiling
at the coupling energy:
\begin{equation}
\frac{\dd}{\dd t}\ln\Meff(t)
\;=\;\frac{\dd \Svn(\sysrho)}{\dd t}
\;\leq\;\frac{2\,\Delta E_{\mathrm{SB}}}{\pi\hbar}
\;=\;\frac{2}{\pi}\,b .
\label{eq:counting-bound}
\end{equation}
The transverse relaxation rate, defined operationally as the initial
decay rate of the transverse coherence
$\mathcal{C}(t)=|\Tr(\sysrho(t)\,\sigma_{+})|$, is bounded by the same
ceiling,
\begin{equation}
\frac{1}{\Ttwo}\;\leq\;\frac{2}{\pi}\,b
\;=\;\frac{2}{\pi}\,\langle\delta\omega^{2}\rangle^{1/2}.
\label{eq:T2-bound}
\end{equation}
\end{theorem}

\begin{proof}
The bath states correlated with distinct spin configurations become
distinguishable no faster than the ceiling \eqref{eq:rate-ceiling} allows,
with $\Delta E=\Delta E_{\mathrm{SB}}$. Each newly distinguishable bath
state adds at most one unit to the participation number of the Schmidt
decomposition, i.e.\ increments $\Meff$ by at most a factor controlled by
$\dd N_{\perp}/\dd t$; taking logarithms, $\dd\ln\Meff/\dd t\leq
2\Delta E_{\mathrm{SB}}/\pi\hbar$, which is \eqref{eq:counting-bound}. For
\eqref{eq:T2-bound}: transverse coherence is the off-diagonal weight of
$\sysrho$, and its decay is bounded below in time (equivalently the rate
bounded above) by the rate at which the bath resolves the spin's phase,
which is exactly the distinguishability rate just bounded. Substituting
$\Delta E_{\mathrm{SB}}=\hbar b$ gives \eqref{eq:T2-bound}.
\end{proof}

\subsection{Saturation reproduces the Anderson--Kubo rate}

The ceiling \eqref{eq:T2-bound} is an inequality. Its physical content is
fixed by asking when it is saturated.

\begin{proposition}[Rigid-lattice saturation]
\label{prop:saturation}
In the slow-bath (rigid-lattice) limit, where the bath correlation time
$\tau_{c}$ is long compared with $1/b$, the ceiling
\eqref{eq:T2-bound} is saturated up to an $O(1)$ geometric factor, and the
relaxation rate is
\begin{equation}
\frac{1}{\Ttwo}\;\simeq\;\langle\delta\omega^{2}\rangle^{1/2}
\;=\;b,
\label{eq:AK}
\end{equation}
the Anderson--Kubo rigid-lattice second-moment rate
\citep{anderson1953,kubo1962}, with no fitted proportionality constant.
\end{proposition}

\begin{proof}
When $\tau_{c}\gg 1/b$ the conditional bath Hamiltonians are effectively
static over the coherence time, so the phase accumulated by the spin is
$\varphi(t)\approx\delta\omega\,t$ with $\delta\omega$ drawn from a
distribution of width $b$. The coherence is the Gaussian average
$\mathcal{C}(t)=\langle e^{i\delta\omega t}\rangle
=\exp(-\tfrac12 b^{2}t^{2})$, whose $1/e$ time is
$\sqrt{2}/b$. Identifying the decay rate gives $1/\Ttwo\simeq b$, which
saturates \eqref{eq:T2-bound} to within the factor $\pi/2$ (the difference
between the Gaussian $1/e$ time and the orthogonality time of the
Margolus--Levitin bound). The motional-narrowing limit $\tau_{c}\ll 1/b$
falls strictly below the ceiling, $1/\Ttwo\simeq b^{2}\tau_{c}\ll b$, as
the standard theory also gives \citep{kubo1962}: the system is then far
from saturating the speed limit, consistent with the inequality.
\end{proof}

\begin{remark}[What has and has not been derived]
\label{rem:ceiling-status}
Proposition~\ref{prop:saturation} reproduces the known rigid-lattice rate
from a speed-limit ceiling, fixing the proportionality constant by
saturation rather than by fitting. It does not, by itself, exceed
Bloch--Redfield theory: in the rigid-lattice limit both give $1/\Ttwo\sim
b$, and in the motional-narrowing limit both give $1/\Ttwo\sim
b^{2}\tau_{c}$. The value added at this layer is conceptual: the rate is
exhibited as the saturation of a universal bound on how fast a bath can
acquire information, not as the output of a particular correlation
function. The genuinely new prediction appears only at the third layer
(Section~\ref{sec:floor}), where the bound is applied to many-body
trajectory counting and yields a coupling-independent floor. The bridge to
that result is the mechanism of the next section, which is required to
know that the count grows monotonically at all.
\end{remark}

% =====================================================================
\section{Mechanism: monotonicity from the partial trace}
\label{sec:mechanism}
% =====================================================================

The ceiling of Section~\ref{sec:ceiling} bounds how fast $\Meff$ can grow.
It does not say that $\Meff$ grows, nor that it grows only forward.
Loschmidt's objection is precisely that a reversible substrate should not
single out a direction. We now show that the count is monotone
non-decreasing under the coarse-grained dynamics, that this monotonicity
is a theorem about the partial trace rather than a postulate, and that the
spin echo realizes the distinction between the reversible and irreversible
parts of dephasing exactly.

\subsection{The reduced dynamics is a CPTP map}

\begin{lemma}[Reduced evolution is completely positive and
trace-preserving]
\label{lem:cptp}
Under Assumptions~\ref{ass:reversible}--\ref{ass:product}, the map
$\Lambda_{t}:\sysrho(0)\mapsto\sysrho(t)$ defined by
\begin{equation}
\Lambda_{t}(\sysrho(0))
=\Tr_{\mathrm{B}}\!\Bigl[U(t)\,
\bigl(\sysrho(0)\otimes\rho_{\mathrm{B}}(0)\bigr)\,U(t)^{\dagger}\Bigr]
\label{eq:lambda}
\end{equation}
is completely positive and trace preserving (CPTP), and admits a
Kraus representation $\Lambda_{t}(\cdot)=\sum_{k}K_{k}(t)\,(\cdot)\,
K_{k}(t)^{\dagger}$ with $\sum_{k}K_{k}^{\dagger}K_{k}=\mathbb{1}$.
\end{lemma}

\begin{proof}
This is the Stinespring/Kraus construction \citep{nielsen2010,breuer2002}.
With an initial product state (Assumption~\ref{ass:product}), let
$\rho_{\mathrm{B}}(0)=\sum_{\mu}p_{\mu}|\mu\rangle\langle\mu|$ in its
eigenbasis. Then
$K_{k}(t)=\sqrt{p_{\mu}}\,\langle\nu|U(t)|\mu\rangle$ with composite index
$k=(\mu,\nu)$ are operators on $\Hilb_{\mathrm{S}}$, unitarity of $U$ gives
$\sum_{k}K_{k}^{\dagger}K_{k}=\mathbb{1}$, and \eqref{eq:lambda} is
recovered. Complete positivity and trace preservation are immediate from
the Kraus form.
\end{proof}

\begin{remark}[Where the arrow enters]
\label{rem:arrow}
Lemma~\ref{lem:cptp} is the entire source of irreversibility. The total
map $\rho\mapsto U\rho U^{\dagger}$ is unitary and reversible (its inverse
is $U^{\dagger}\,\cdot\,U$, equally physical). The reduced map
$\Lambda_{t}$ is in general \emph{not} unitarily invertible: the partial
trace $\Tr_{\mathrm{B}}$ discards the spin--bath correlations, and once
discarded they cannot be recovered by any operation on the spins alone.
The arrow of time in this framework is neither postulated nor smuggled
into the definition of a count; it is the non-invertibility of the partial
trace, the one operation in the construction that fails to commute with
time reversal.
\end{remark}

\subsection{Monotonicity of the count}

\begin{theorem}[Monotone partition count]
\label{thm:monotone}
Let the bath be initially in a maximally mixed (infinite-temperature)
state, $\rho_{\mathrm{B}}(0)\propto\mathbb{1}$, as is standard for nuclear
spin baths at ordinary temperatures \citep{abragam1961}. Then the partition
count is non-decreasing along the coarse-grained evolution:
\begin{equation}
\Meff(t_{2})\;\geq\;\Meff(t_{1})\qquad\text{for all }t_{2}\geq t_{1},
\label{eq:monotone}
\end{equation}
equivalently $\Svn(\sysrho(t_{2}))\geq\Svn(\sysrho(t_{1}))$.
\end{theorem}

\begin{proof}
For a maximally mixed bath, $\Lambda_{t}$ is a \emph{unital} CPTP map
(it preserves the maximally mixed system state, since the global dynamics
cannot concentrate probability when the bath supplies no bias). Unital
CPTP maps are entropy non-decreasing: by the monotonicity of relative
entropy under CPTP maps \citep{lindblad1975} applied to the maximally
mixed reference $\pi=\mathbb{1}/d_{\mathrm{S}}$,
\begin{equation}
S(\Lambda_{t}\sysrho\,\|\,\pi)\leq S(\sysrho\,\|\,\pi),
\end{equation}
and since $S(\sigma\|\pi)=\ln d_{\mathrm{S}}-\Svn(\sigma)$ for unital
$\Lambda_{t}$, this rearranges to
$\Svn(\Lambda_{t}\sysrho)\geq\Svn(\sysrho)$. Composing maps over
$[t_{1},t_{2}]$ (the family $\Lambda_{t}$ is a one-parameter CP-divisible
evolution in the Markovian regime; for the non-Markovian case see
Remark~\ref{rem:nonmarkov}) gives \eqref{eq:monotone}. Exponentiating
yields the statement for $\Meff$.
\end{proof}

\begin{corollary}[Reversible substrate, irreversible count]
\label{cor:reversible-irreversible}
The total von Neumann entropy $\Svn(\rho(t))$ is constant (unitary
invariance), while the reduced entropy $\Svn(\sysrho(t))$ is
non-decreasing. The conserved global entropy is redistributed into
spin--bath correlations; the monotone reduced entropy is the count. No
contradiction with Assumption~\ref{ass:reversible} arises, because the two
entropies refer to different states.
\end{corollary}

\begin{remark}[Non-Markovian baths and the softening of strict
monotonicity]
\label{rem:nonmarkov}
For a genuinely non-Markovian bath the family $\{\Lambda_{t}\}$ need not be
CP-divisible, and $\Svn(\sysrho(t))$ can transiently \emph{decrease}---this
is information backflow, the microscopic basis of partial echo revivals. We
therefore state the honest form of the result: $\Meff(t)$ is monotone
non-decreasing on timescales long compared with the bath correlation time
$\tau_{c}$, and admits bounded transient dips on shorter scales whose
integrated effect is recovered by refocusing. Strict monotonicity holds in
the Markovian (coarse-grained) limit; controlled non-monotonicity on the
scale $\tau_{c}$ is not a defect of the theory but its prediction for when
echoes can partially revive. This ties the monotonicity of the count
directly to the refocusability structure of the next subsection.
\end{remark}

\subsection{The spin echo separates the two parts exactly}

We now make precise the claim that the Hahn echo recovers exactly the part
of dephasing that did not increment the count.

\begin{theorem}[Echo refocuses the product-state part only]
\label{thm:echo-split}
Write the dephasing of a spin as the sum of a static contribution (from a
fixed offset $\Delta\omega$ that is a property of the spin's location, not
of the bath) and a fluctuating contribution (from $H_{\mathrm{SB}}$). Then:
\begin{enumerate}[leftmargin=1.5em,label=(\roman*),itemsep=0.2em]
\item The static contribution keeps the spin--bath state a product:
$\Meff$ is unchanged, $\Svn(\sysrho)=0$ throughout. A $180^{\circ}$ pulse
at $\tau$ maps the accumulated phase $\varphi\mapsto-\varphi$ and the
ensemble refocuses exactly at $2\tau$. The dephasing is reversible because
no count was incremented.
\item The fluctuating contribution entangles the spin with the bath:
$\Meff$ grows, $\Svn(\sysrho)>0$. A $180^{\circ}$ pulse acts on the spin
alone and cannot disentangle it from the bath; the residual coherence at
$2\tau$ is reduced by $\exp[-\Delta\Svn]$, where $\Delta\Svn$ is the
entropy committed in $[0,2\tau]$. The dephasing is irreversible because the
count was incremented.
\end{enumerate}
\end{theorem}

\begin{proof}
(i) A static offset commutes with the partial trace: the spin phase is a
$c$-number function of a fixed parameter, so $U(t)$ restricted to that
contribution acts as a local unitary on $\Hilb_{\mathrm{S}}$, leaving the
state a product and $\Svn(\sysrho)=0$. A local unitary ($180^{\circ}$
pulse) inverts the phase, and continued evolution under the same static
offset cancels it at $2\tau$, independently of the offset value. (ii) The
fluctuating coupling $H_{\mathrm{SB}}$ generates genuine spin--bath
entanglement (the conditional bath states correlated with distinct spin
phases become non-identical, hence non-orthogonal-to-orthogonal as in
Theorem~\ref{thm:counting-bound}), so $\Svn(\sysrho)>0$. By monogamy of
entanglement and the locality of the pulse \citep{coffman2000}, no operation
on $\Hilb_{\mathrm{S}}$ alone can reduce $\Svn(\sysrho)$; the committed
correlation is inaccessible to refocusing, and the echo amplitude carries
the factor $e^{-\Delta\Svn}$. Identifying $\Delta\Svn=2\tau/\Ttwo$ recovers
the standard echo attenuation $e^{-2\tau/\Ttwo}$.
\end{proof}

\begin{remark}[The qualitative becomes quantitative]
\label{rem:quantitative}
Theorem~\ref{thm:echo-split} converts the familiar statement ``$\Tstar$ is
refocusable and $\Ttwo$ is not'' into a statement about a computable
quantity: refocusable $\Leftrightarrow$ $\Svn(\sysrho)=0$ (product state),
non-refocusable $\Leftrightarrow$ $\Svn(\sysrho)>0$ (entangled state), and
the echo deficit is exactly $e^{-\Delta\Svn}$ with $\Delta\Svn$ the
entropy committed to the bath. The committed entropy is not a free
parameter: its rate is bounded by Theorem~\ref{thm:counting-bound} and, in
the rigid-lattice limit, saturates to $b$ by
Proposition~\ref{prop:saturation}. The identification
$\Delta\Svn=2\tau/\Ttwo$ is therefore a derived relation, $1/\Ttwo$ being
the committed-entropy rate.
\end{remark}

% =====================================================================
\section{Headline: a perturbation-independent Loschmidt-echo floor}
\label{sec:floor}
% =====================================================================

We now reach the result that distinguishes the counting framework from the
correlation-function account. The Loschmidt (polarization) echo reverses
the \emph{many-body} interaction, $H\to-H+\Sigma$, where $\Sigma$ collects
the residual un-reversed couplings. The empirical discovery
\citep{levstein1998,pastawski2000,jalabert2001} is that, in the
strongly-coupled regime, the echo-fidelity decay rate saturates at a value
\emph{independent of $\|\Sigma\|$}: improving the reversal does not slow the
decay. We derive this saturation, and its rate, from trajectory counting.

\subsection{Counting many-body trajectories}

\begin{definition}[Channel multiplicity]
\label{def:channels}
At each coherent step the spin system can commit which-state information
through any of $d$ independent coupling channels (e.g.\ the distinct
neighbours of a spin in a dipolar network, or the independent decay
modes available to the local cluster). $d$ is a structural property of the
coupling graph, not of the coupling strength.
\end{definition}

\begin{theorem}[Trajectory-count growth]
\label{thm:trajectory-count}
The number of mutually distinguishable labelled commitment-trajectories the
system can traverse in $n$ coherent steps over $d$ channels is
\begin{equation}
T(n,d)=\sum_{k=1}^{n}\binom{n-1}{k-1}d^{k}=d\,(d+1)^{n-1}.
\label{eq:Tnd}
\end{equation}
Its logarithmic growth rate per step is
\begin{equation}
\lambda_{\mathrm{count}}
:=\lim_{n\to\infty}\frac{1}{n}\ln T(n,d)=\ln(d+1),
\label{eq:lambda}
\end{equation}
independent of any coupling strength.
\end{theorem}

\begin{proof}
A commitment-trajectory of $n$ steps is an ordered composition of $n$ into
$k$ parts (the run-lengths between channel switches), each part labelled by
one of $d$ channels. Compositions of $n$ into exactly $k$ positive parts
number $\binom{n-1}{k-1}$ (stars and bars); each part carries one of $d$
labels, contributing $d^{k}$. Summing,
$T(n,d)=\sum_{k=1}^{n}\binom{n-1}{k-1}d^{k}
=d\sum_{j=0}^{n-1}\binom{n-1}{j}d^{j}=d(1+d)^{n-1}$ by the binomial
theorem. Then $\tfrac1n\ln T=\tfrac1n[\ln d+(n-1)\ln(d+1)]\to\ln(d+1)$.
The result depends only on the combinatorial multiplicity $d$, with no
reference to $H_{\mathrm{SB}}$'s magnitude or to $\Sigma$.
\end{proof}

\subsection{The floor rate}

\begin{theorem}[Perturbation-independent echo floor]
\label{thm:floor}
Let $\omega_{\mathrm{loc}}$ be the local precession (commitment) frequency,
so that one coherent step occupies time $2\pi/\omega_{\mathrm{loc}}$. In
the strongly-coupled regime the Loschmidt-echo fidelity decays no slower
than the rate at which distinguishable trajectories proliferate,
\begin{equation}
\boxed{\;\frac{1}{\tau_{\mathrm{floor}}}
=\frac{\omega_{\mathrm{loc}}}{2\pi}\,\ln(d+1)\;}
\label{eq:floor}
\end{equation}
and this rate is independent of the residual perturbation $\|\Sigma\|$ used
to engineer the reversal.
\end{theorem}

\begin{proof}
The fidelity $F(t)=|\langle\psi_{0}|e^{iH't/\hbar}e^{-iHt/\hbar}|\psi_{0}
\rangle|^{2}$ is the squared overlap of the forward-evolved state with the
imperfectly back-evolved state. By Theorem~\ref{thm:echo-split}(ii), the
recoverable part of the overlap is limited by the committed entropy
$\Svn(\sysrho)$, which grows as the system explores distinguishable
commitment-trajectories. After $n$ steps the number of such trajectories
is $T(n,d)$, so $\Svn\geq\ln(\text{effective participation})$ grows at rate
$\lambda_{\mathrm{count}}=\ln(d+1)$ per step
(Theorem~\ref{thm:trajectory-count}). Converting steps to time through
$\omega_{\mathrm{loc}}/2\pi$ steps per unit time gives the decay rate
\eqref{eq:floor}. Independence of $\|\Sigma\|$ is inherited from
\eqref{eq:lambda}: the proliferation rate is set by the channel
multiplicity $d$ alone, so reducing the residual coupling (improving the
reversal) does not reduce it. The reversal can recover the configurational
overlap but not the trajectory-counting entropy, exactly as the Hahn echo
recovers $\Tstar$ but not $\Ttwo$.
\end{proof}

\begin{corollary}[Numerical value for a three-dimensional network]
\label{cor:numerics}
For a spin embedded in a three-dimensional dipolar network with effective
channel multiplicity $d=3$ (three independent angular commitment
directions), $\ln(d+1)=\ln 4\approx 1.386$, and
\begin{equation}
\frac{1}{\tau_{\mathrm{floor}}}\approx 0.221\,\omega_{\mathrm{loc}}.
\label{eq:numeric-floor}
\end{equation}
Identifying $\omega_{\mathrm{loc}}$ with the second moment of the local
field, $\omega_{\mathrm{loc}}=\langle\delta\omega^{2}\rangle^{1/2}=b$, gives
$1/\tau_{\mathrm{floor}}\approx 0.22\,b$: the floor is a fixed fraction of
the rigid-lattice linewidth, set by the network dimensionality.
\end{corollary}

\subsection{Why this is not the semiclassical Lyapunov result}

\begin{remark}[Point of departure from the Lyapunov account]
\label{rem:lyapunov}
The semiclassical theory of the Loschmidt echo in classically chaotic
systems \citep{jalabert2001,gorin2006} also predicts a
perturbation-independent decay, identifying the saturated rate with a
classical Lyapunov exponent $\lambda_{\mathrm{Lyap}}$ of the underlying
dynamics. Equations~\eqref{eq:lambda} and \eqref{eq:floor} agree with that
prediction in form---a coupling-independent exponential rate---but differ
in origin and, crucially, in domain. The counting rate $\ln(d+1)$ is a
property of the \emph{coupling-graph multiplicity}, not of dynamical
instability. It is therefore nonzero for systems whose classical limit is
integrable or non-chaotic, where $\lambda_{\mathrm{Lyap}}=0$ and the
semiclassical theory predicts no floor. This is a sharp, falsifiable point
of departure: a discrete, integrable, or few-body spin cluster with no
positive Lyapunov exponent should, on the counting account, still exhibit a
Loschmidt-echo floor at rate $(\omega_{\mathrm{loc}}/2\pi)\ln(d+1)$ fixed
by its connectivity $d$, whereas the Lyapunov account predicts the floor
recedes to zero as the reversal is perfected.
\end{remark}

\begin{remark}[Consistency with the chaotic regime]
\label{rem:consistency}
When the spin network \emph{is} chaotic, the connectivity-set rate
$\ln(d+1)$ and the Lyapunov rate $\lambda_{\mathrm{Lyap}}$ need not
coincide; the observed floor is then the slower-saturating of the two
mechanisms (the system decoheres at whichever rate first exhausts
recoverable overlap). In dense dipolar solids the two are typically of the
same order---$\ln 4\approx 1.4$ versus Lyapunov exponents of order unity in
units of the local coupling---which is why the existing data
\citep{levstein1998,pastawski2000} are consistent with either reading. The
discriminating experiments are therefore in the \emph{non-chaotic} regime
(Section~\ref{sec:experiments}).
\end{remark}

% =====================================================================
\section{Experimental tests}
\label{sec:experiments}
% =====================================================================

The framework makes three quantitative claims, each independently
testable.

\subsection{Test 1: rigid-lattice rate as a saturated ceiling}

Proposition~\ref{prop:saturation} predicts $1/\Ttwo\simeq b$ in the
rigid-lattice limit with no fitted constant. This is already confirmed by
the standard second-moment analysis of solid-state lineshapes
\citep{abragam1961,anderson1953}; the contribution here is interpretive
(the rate is the saturated Margolus--Levitin ceiling), so the test is one
of \emph{consistency}: in every system where $\tau_{c}\gg 1/b$, the measured
$1/\Ttwo$ should equal the van Vleck second moment $b$ to within the
geometric $O(1)$ factor of Proposition~\ref{prop:saturation}, and should
fall strictly below $b$ as $\tau_{c}$ decreases into motional narrowing.

\subsection{Test 2: the floor's connectivity dependence}

Corollary~\ref{cor:numerics} predicts $1/\tau_{\mathrm{floor}}=
(\omega_{\mathrm{loc}}/2\pi)\ln(d+1)$. Two consequences are directly
measurable. First, the floor rate should scale as $\ln(d+1)$ as the
coordination number of the spin network is varied (e.g.\ by isotopic
dilution, which reduces the number of active dipolar partners $d$): a plot
of $\tau_{\mathrm{floor}}^{-1}/\omega_{\mathrm{loc}}$ against $\ln(d+1)$
should be a straight line of unit slope through the origin. Second, the
floor should be \emph{independent} of the quality of the magic-echo or
time-reversal pulse sequence once the reversal is good enough to enter the
saturated regime---reproducing the observed environment-independence
\citep{pastawski2000} but now with a predicted magnitude.

\subsection{Test 3: a floor in a non-chaotic system (the decisive test)}

Remark~\ref{rem:lyapunov} gives the experiment that separates the counting
account from the Lyapunov account. Prepare a small, well-characterized,
\emph{non-chaotic} or integrable spin cluster (for instance a linear chain
with nearest-neighbour-only XY coupling, or a few-spin star geometry) whose
classical limit has $\lambda_{\mathrm{Lyap}}=0$. Perform a polarization-echo
reversal of increasing quality.
\begin{itemize}[leftmargin=1.4em,itemsep=0.2em]
\item \emph{Lyapunov prediction:} the echo decay rate $\to 0$ as
$\|\Sigma\|\to 0$; the floor recedes.
\item \emph{Counting prediction:} the echo decay rate saturates at
$(\omega_{\mathrm{loc}}/2\pi)\ln(d+1)$ with $d$ the cluster's connectivity;
the floor remains.
\end{itemize}
Observation of a nonzero, connectivity-set floor in a provably non-chaotic
cluster would confirm the counting mechanism and exclude the pure Lyapunov
account; observation of a receding floor would refute the counting
mechanism at this level. Either outcome is informative.

% =====================================================================
\section{Discussion}
\label{sec:discussion}
% =====================================================================

\subsection{What resolves the Loschmidt objection, and what does not}

The Loschmidt objection is that a time-reversal-invariant substrate cannot
single out a temporal direction. The resolution offered here is precise
about where the asymmetry lives. It is not in the substrate---%
Assumption~\ref{ass:reversible} keeps the total Hamiltonian reversible and
Corollary~\ref{cor:reversible-irreversible} keeps the total entropy
constant. It is not in a definition---the count $\Meff$ is the effective
Schmidt rank of a genuine quantum state, and its monotonicity is a theorem
(Theorem~\ref{thm:monotone}), not a stipulation. The asymmetry enters at
exactly one place: the partial trace $\Tr_{\mathrm{B}}$
(Remark~\ref{rem:arrow}), the operation that discards spin--bath
correlations and that alone fails to commute with time reversal. The second
law, in this reading, is the statement that the reduced count cannot
decrease because the discarded correlations cannot be recalled by local
operations.

This is a conventional resolution in its physics---it is the open-systems,
decoherence-theoretic account \citep{zurek2003,breuer2002}---but the paper
makes two contributions beyond restating it. First, it fixes the
\emph{rate} of the count's growth by a universal speed-limit ceiling
(Theorem~\ref{thm:counting-bound}) rather than by a model-specific
correlation function, and shows the relaxation rate is its saturation
(Proposition~\ref{prop:saturation}). Second, it derives a
\emph{coupling-independent echo floor} (Theorem~\ref{thm:floor}) from pure
trajectory counting, with a connectivity-set magnitude and a falsifiable
prediction in the non-chaotic regime (Test~3).

\subsection{Honest limitations}

We list the load-bearing soft spots so they can be attacked directly.
\begin{enumerate}[leftmargin=1.5em,itemsep=0.2em]
\item \emph{The step-to-time conversion.} Equation~\eqref{eq:floor}
identifies the coherent-step rate with $\omega_{\mathrm{loc}}/2\pi$ and,
in Corollary~\ref{cor:numerics}, $\omega_{\mathrm{loc}}$ with the
second-moment frequency $b$. This identification is physically motivated
but not derived from first principles; a factor-of-order-unity geometric
correction is possible and must be pinned down against data before the
numerical coefficient $0.22$ is treated as exact.
\item \emph{The channel multiplicity $d$.} The value $d=3$
(Corollary~\ref{cor:numerics}) is the natural choice for an isotropic
three-dimensional dipolar network (three independent angular commitment
directions), but the effective $d$ for a given lattice is a modelling
choice that the connectivity-scaling test (Test~2) is designed to
calibrate, not assume.
\item \emph{Strict versus coarse-grained monotonicity.} Theorem~%
\ref{thm:monotone} gives strict monotonicity only in the Markovian limit;
non-Markovian baths permit bounded transient decreases of the count
(Remark~\ref{rem:nonmarkov}). The framework predicts these as partial echo
revivals rather than excluding them, but the claim ``the count never
decreases'' must be qualified to ``the count never decreases on scales long
compared with $\tau_{c}$.''
\item \emph{Saturation, not equality.} The relaxation rate equals the
ceiling only at saturation (Proposition~\ref{prop:saturation}); away from
the rigid-lattice limit the ceiling is a strict inequality and the actual
rate is set by the bath spectral density, as in the standard theory. The
counting framework bounds and, in one limit, reproduces the standard
result; it does not replace it everywhere.
\end{enumerate}

\subsection{Outlook}

The decisive next step is Test~3: a Loschmidt-echo measurement on a
provably non-chaotic spin cluster. The counting framework predicts a
nonzero floor fixed by connectivity where the semiclassical Lyapunov
account predicts none. If the floor is observed with the predicted
connectivity dependence, the relaxation rate, the echo deficit, and the
many-body fidelity floor are unified as three manifestations of a single
counting bound saturating a quantum speed limit. If it is not, the
trajectory-counting layer (Section~\ref{sec:floor}) is falsified while the
ceiling (Section~\ref{sec:ceiling}) and the monotonicity mechanism
(Section~\ref{sec:mechanism}) stand independently.

% =====================================================================
\section{Conclusion}
\label{sec:conclusion}
% =====================================================================

We have given an independent, information-theoretic account of transverse
relaxation and echo attenuation built on three layers. A counting ceiling
derived from the Margolus--Levitin quantum speed limit bounds the rate at
which a bath can acquire which-spin information, and its saturation
reproduces the Anderson--Kubo rigid-lattice relaxation rate
$1/\Ttwo\simeq\langle\delta\omega^{2}\rangle^{1/2}$ with no fitted constant.
The monotonicity of the underlying count is derived from the
non-invertibility of the partial trace---the arrow of time entering through
open-system coarse-graining of a strictly reversible substrate, not through
any postulate---and the Hahn echo is shown to separate the reversible
(product-state) and irreversible (entanglement-committed) parts of dephasing
exactly. Counting the distinguishable many-body commitment-trajectories then
yields a perturbation-independent Loschmidt-echo floor
$1/\tau_{\mathrm{floor}}=(\omega_{\mathrm{loc}}/2\pi)\ln(d+1)$, reproducing
the observed environment-independence of polarization-echo decay and
predicting---unlike the semiclassical Lyapunov account---a nonzero floor for
non-chaotic spin systems. The framework is quantitative throughout, makes a
falsifiable prediction that departs from the orthodox account in a specific
regime, and locates the irreversibility of the spin echo in a single,
identifiable place.

\bibliographystyle{plainnat}
\bibliography{references}

\end{document}